\subsubsection{Variáveis}

Variáveis podem ser entendidas como objetos com nomes personalizados, 
capazes de armazenar valores que serão utilizados pelo programa durante 
sua execução. A declaração de uma variável geralmente segue o seguinte padrão:

\begin{textbox}
    <tipo da variável> <nome da variável>;
\end{textbox}

Os nomes das variáveis podem conter letras e números, além de poderem 
ser iniciados com um sublinhado. No entanto, algumas regras devem ser 
seguidas para declarar uma variável sem problemas.

Na linguagem C, algumas palavras já são reservadas para representar 
diretivas e, portanto, não podem ser usadas como nomes de variáveis. 
Abaixo, segue uma tabela das palavras reservadas:

\begin{table}[h]
  \centering
  \begin{tabular}{llll}
    auto      & break     & case      & char      \\
    const     & continue  & default   & do        \\
    double    & else      & enum      & extern    \\
    float     & for       & goto      & if        \\
    int       & long      & register  & return    \\
    short     & signed    & sizeof    & static    \\
    struct    & switch    & typedef   & union     \\
    unsigned  & void      & volatile  & while     \\
  \end{tabular}
\end{table}

Você não precisa necessariamente memorizar todas essas diretivas neste 
momento; isso é algo que você absorverá e memorizará com bastante prática e 
um pouco de tempo.

\subsubsection{Tipos de dados} \label{tipos-de-dados}

Um aspecto que ainda não abordamos é a 'tipagem' de variáveis. Embora cada 
variável tenha um nome definido, o computador não sabe exatamente o que há 
dentro de cada uma delas. Portanto, cabe ao usuário fazer essa definição.

A princípio, a linguagem C fornece quatro tipos básicos: \textbf{char}, 
\textbf{int}, \textbf{float} e \textbf{double}. Além disso, há modificadores 
como \textbf{signed}, \textbf{unsigned}, \textbf{short} e \textbf{long}, que 
podem ser usados em conjunto com cada tipo.

Novamente, não é necessário memorizar tudo isso agora, mas por curiosidade, na 
próxima página terá uma lista descrevendo os tipos e modificadores, juntamente 
com algumas informações adicionais:

\newpage

\begin{center}
    \centering
    \begin{tabular}{|c|p{4cm}|p{2cm}|c|p{4cm}|}
        \hline
        \textbf{Tipo} & \textbf{Descrição} & \textbf{Tamanho (bits)} & \textbf{Formatador} & \textbf{Alcance} \\
        \hline
        char & Representa um byte de dados. Pode ser utilizado para caracteres ou valores numéricos pequenos. Pode ser \textit{signed} ou \textit{unsigned}. & 8 & \%c (ou \%d) & -128 a 127 (signed) ou 0 a 255 (unsigned) \\
        \hline
        int & Tipo inteiro padrão, utilizado para operações aritméticas gerais. & 32 & \%d (ou \%i) & aproximadamente $-2^{31}$ a $2^{31}-1$ \\
        \hline
        float & Representa números de ponto flutuante de precisão simples. & 32 & \%f & $\approx 1.17 \times 10^{-38}$ a $3.40 \times 10^{38}$ \\
        \hline
        double & Representa números de ponto flutuante de dupla precisão. & 64 & \%lf & $\approx 2.22 \times 10^{-308}$ a $1.79 \times 10^{308}$ \\
        \hline
    \end{tabular}
\end{center}

Além dos tipos de dados descritos acima, segue abaixo, uma descrição mais detalhada do que cada um dos 
modificadores de tipo fazem:

\begin{center}
    \begin{tabular}{|c|p{11cm}|}
        \hline
        \textbf{Modificador} & \textbf{Descrição} \\
        \hline
        signed & Indica que o tipo suporta valores negativos (com sinal). Geralmente é o padrão para tipos inteiros. \\
        \hline
        unsigned & Indica que o tipo armazena apenas valores não negativos, dobrando o limite superior positivo. \\
        \hline
        short & Reduz o tamanho de um inteiro (ex: \texttt{short int}). Garante no mínimo 16 bits. \\
        \hline
        long & Aumenta o tamanho de um inteiro (ex: \texttt{long int}). Garante no mínimo 32 bits. \\
        \hline
        long long & Inteiro de maior capacidade, garantindo no mínimo 64 bits. \\
        \hline
    \end{tabular}
\end{center}

\clearpage

Agora, saindo um pouco da teoria e partindo para a prática, vamos aplicar os 
conceitos aprendidos em um programa simples.

\begin{excode}
    #include <stdio.h>

    int main(void) {
        // Declara variáveis para guardar valores
        // necessários para realizar os cálculos
        float x;
        float y;
        float area;

        // Pede ao usuário por valores de largura e altura
        printf("Digite a altura do retângulo: ");
        scanf("%f", &x);

        printf("Digite a comprimento do retângulo: ");
        scanf("%f", &y);

        // Calcula a área do triângulo
        area = x * y;

        // Escreve o resultado no terminal
        printf("A área do retângulo é: %f\n", area);

        return 0;
    }
\end{excode}

O código acima funciona da seguinte forma: ao ser iniciado, o usuário deve 
digitar no terminal os valores da altura e comprimento de um retângulo 
imaginário. Em seguida, o programa calcula automaticamente a área do 
retângulo e exibe o resultado.

\begin{excode}
        float x;
        float y;
        float area;
\end{excode}

No trecho abaixo, declaramos três variáveis diferentes: \texttt{x} para a 
altura, \texttt{y} para o comprimento e \texttt{area} para o resultado do 
cálculo. Todas as três variáveis são do tipo \textbf{float}, pois nada impede 
que as medidas do retângulo sejam números fracionados.

\begin{excode}
        printf("Digite a altura do retângulo: ");
        scanf("%f", &x);

        printf("Digite a comprimento do retângulo: ");
        scanf("%f", &y);
\end{excode}

Aqui utilizamos a função \texttt{printf} para exibir texto na tela, assim 
como já feito antes, e em seguida, usamos uma nova função chamada 
\texttt{scanf} para receber informações do usuário pelo terminal.

Dentro de scanf, passamos dois parâmetros: o primeiro (\texttt{\%f}) 
é o formatador, que especifica o tipo de valor esperado para ser recebido, 
neste caso, uma \textbf{float}. O segundo (\texttt{\&x}) é onde queremos 
armazenar o valor recebido; a variável vem com o sufixo \& porque estamos 
passando o endereço da variável.

\begin{postscript}
    O endereço de uma variável é um tópico que será abordado com 
    maior foco no capítulo \ref{o-que-sao-enderecos}, porém, no momento você pode assumir que ele 
    diz respeito a localização da variável na memória do computador.
\end{postscript}

\begin{excode}
        area = x * y;
\end{excode}

Após receber os valores de \(x\) e \(y\) usando a função \texttt{scanf}, 
multiplicamos as variáveis usando o operador de multiplicação e armazenamos 
o resultado em \texttt{area} usando o operador de atribuição (\texttt{=}).

\begin{excode}
        printf("A área do retângulo é: %f\n", area);
        return 0;
\end{excode}

Finalmente, exibimos o valor na tela usando a função \texttt{printf} e 
retornamos o valor 0 para indicar que o programa foi encerrado com sucesso. 
É importante lembrar que, ao exibir uma variável do tipo \textbf{float}, é 
necessário utilizar o formatador \texttt{\%f}.

Entendendo bem o código acima, é possível praticar os conceitos de operadores 
e variáveis sem muitas dificuldades.

\subsubsection{Exercícios}

\textbf{É de extrema importância que você realize os exercícios aqui 
elaborados para fixar o conteúdo.}

\begin{enumerate}
    \item{Modifique a lógica do programa acima para calcular a área de um 
        triângulo.}
    \item{Crie uma calculadora simples capaz de receber 2 números e realizar 
        alguma das 4 operações básicas.}
    \item{Crie um conversor de unidade de metro para centímetros.}
\end{enumerate}
